% Owner: B
We use the course definitions of centrality \cite{quercini2026}. Degree counts
neighbours; strength sums incident edge weights. We compute both and hop
betweenness on each graph. We use PMI weights for ingredients and Jaccard
weights for countries.

Betweenness sums, over pairs of other nodes, the share of their shortest
paths that pass through the node considered. Our primary measure of brokers,
meaning ingredients that lie on many shortest paths, is hop betweenness.
We give every edge the same length. NetworkX normalises by the number of
pairs of other nodes: the value is the average share of their shortest paths
passing through the node. Country betweenness uses the whole graph; pairs
in different components contribute nothing.

For sensitivity, we compute betweenness with edge distance $1/\mathrm{PMI}$
and a second conversion, $d=1/(1+\mathrm{PMI})$. Any positive decreasing
function of PMI can define edge lengths, making stronger associations shorter.
The course gives no rule for choosing one. Here PMI ranges from $0.000636$
to $6.643$, so $1/\mathrm{PMI}$ gives extreme distances up to about $1{,}573$.
The second conversion gives distances of $0.131$--$0.9994$. It is still one
choice among many: agreement tests robustness to these two conversions only.
Hop betweenness remains our primary measure.

We also subtract each ingredient's hop-betweenness rank from its degree rank.
A positive rank gap means it ranks higher in betweenness than in degree.
We compare these two taught measures to find betweenness beyond degree,
restricting discussion to the hop top $50$ to avoid gaps among near-zero values.

To compare degree and the betweenness measures with ingredient frequency,
the number of recipes containing the ingredient, we use Spearman
correlation, which compares ranks. We read these correlations descriptively
because frequency comes from the same recipes as the graph. They do not
establish a causal relationship.
